
\def\PUDcodigo{ACE013}

\if\COMPILAR0
    \def\PUDcodacad{04507.9}
\else
    \def\PUDcodacad{04506.9}
\fi

\def\PUDdisciplina{Probabilidade e Estatística}

%NOVOS ---------------------------------------------------
\def\PUDhorasTeoria{80}
\def\PUDhorasPratica{0}

\def\PUDinterECA{---}

\def\PUDatuaisECA{---}

\def\PUDrecursos{Projetor multimídia; Quadro branco e pincel.}

%----------------------------------------------------------

\def\PUDsemestre{S2}

\def\PUDhoras{80}

\def\PUDcreditos{4}

\def\PUDprerequisito{ \hyperref[PUD:ACE003]{ACE003} }

\if\COMPILAR0
    \def\PUDpreacad{ \hyperref[PUD:ACE003]{Cálculo 1 (04507.2)} }
\else
    \def\PUDpreacad{ \hyperref[PUD:ACE003]{Cálculo 1 (04506.2)} }
\fi

\def\PUDobjetivos{Compreender os conceitos teóricos de Probabilidade e Estatística, através dos seguintes objetivos específicos: Conceituar os elementos básicos da estatística e identificar as etapas de um trabalho estatístico; Caracterizar os tipos de coletas de dados, distinguindo os tipos de variáveis;  Sintetizar os dados através de tabelas, analisando-os por meio dos gráficos; 
Conhecer as aplicações, calculando e interpretando as medidas de posição;  Conhecer as aplicações, calculando e interpretando as medidas de dispersão;  Conhecer as aplicações, calculando e interpretando as medidas de associação;  Avaliar a previsibilidade de dados estatísticos;  Conhecer as premissas que permitem o emprego dos modelos probabilísticos discretos;  Conhecer as premissas que permitem o emprego dos modelos probabilísticos contínuos;  Identificar as relações existentes entre a população e as amostras extraídas;  Realizar inferências acerca de uma população baseada nos dados amostrais.}

\def\PUDementa{Conceitos fundamentais da estatística;  Estudo dos dados estatísticos;  Representação gráfica e tabular;  Medidas de posição;  Medidas de dispersão;  Medidas de associação;  Noções de probabilidade;   Distribuições discretas de probabilidade;  Distribuições contínuas de probabilidade;  Teoria da amostragem;  Estatística indutiva.}

\def\PUDconteudo{ &  Conceitos fundamentais da estatística: Estatística.  População e amostra.  Variável.  Fenômeno determinístico x fenômeno aleatório.  Estatística descritiva e indutiva. 
Parâmetro.  Fases de um trabalho estatístico.  \\ 
 &  Estudo dos dados estatísticos: Variáveis qualitativas: nominal e ordinal.  Variáveis quantitativas: discreta e continua. 
Variáveis descritas em escala nominal, ordinal, intervalar e razões.  Tipos de coletas de dados (periódica, contínua e ocasional).  \\ 
 &  Representação gráfica e tabular: Apresentação dos dados através das séries estatísticas, envolvendo de uma a três variáveis.  Identificação do uso, construção e interpretação dos gráficos em coluna ou barra, linha, pizza, polar.  Sintetização dos dados de uma variável quantitativa em uma tabela de distribuição de frequência.  Cálculo e interpretação das frequências relativas e acumuladas.  Construção e interpretação dos gráficos de segmentos de reta, histograma e polígono.  \\ 
 &  Medidas de posição: Moda.  Mediana.  Média e suas propriedades.  Relações entre moda, mediana e média.  \\ 
 &  Medidas de dispersão: Conceito de dispersão.  Dispersão absoluta (variância e desvio padrão).  Propriedades da variância.  Dispersão relativa (coeficiente de variação e escore reduzido).  \\ 
 &  Medidas de associação: Associação para variável Qualitativa (Teste de Independência, Coeficiente de Contingência).  Associação para variável Quantitativa (Diagrama de Dispersão, Coeficiente de Correlação, Regressão linear, transformações e múltipla, Coeficiente de determinação).  \\ 
 &  Noções de probabilidade: Avaliação da previsibilidade de fenômenos aleatórios.  Distinção entre probabilidade a posteriori e probabilidade a priori.  Teorema da soma, probabilidade condicional e teorema do produto.  Eventos mutuamente exclusivos e independentes.  \\ 
 &  Distribuições discretas de probabilidade: Função densidade de probabilidade.  Representação gráfica.  Distribuição binomial.  Distribuição de Poisson.  \\ 
 &  Distribuições contínuas de probabilidade: Função densidade de probabilidade.  Distribuição normal.  \\ 
 &  Teoria da amostragem: Distribuição amostral das médias.  Distribuição amostral das proporções.  Teorema do limite central.  \\ 
 &  Estatística indutiva: Estimativa para média e diferença de média (para pequenas e grandes amostras).  Estimativa para proporção e diferença de proporção. }

\def\PUDmetodo{Aulas expositórias que podem ser teóricas e/ou práticas, onde as práticas no laboratório serão marcadas no decorrer da disciplina.}

\def\PUDavalia{A avaliação será feita com aplicação de provas teóricas e/ou práticas, além da possibilidade de inclusão de trabalhos e seminários no decorrer da disciplina.}

\def\PUDbibbasica{ & \href{www.biblioteca.ifce.edu.br/index.asp?codigo_sophia=23974}{Magalhães}, M.  N.  Noções de Probabilidade e Estatística.  7º ed.  São Paulo, SP : Universidade de São Paulo, 2010.  408p.  ISBN: 9788531406775. 
\\ 
 & \href{www.biblioteca.ifce.edu.br/index.asp?codigo_sophia=24369}{Crespo}, A.  A.  Estatística Fácil.  19º ed.  São Paulo, SP: Saraiva, 2009.  218p.  ISBN: 9788502081062. 
\\ 
 & \href{www.biblioteca.ifce.edu.br/index.asp?codigo_sophia=24041}{Mucelin}, C.  A.  Estatística.  Curitiba, PR: Livro Técnico, 2010.  120p.  ISBN: 9788563687081. }

\def\PUDbibcomplementar{ & \href{www.biblioteca.ifce.edu.br/index.asp?codigo_sophia=44570}{Albuquerque}, José Paulo de Almeida e.  Probabilidade, variáveis aleatórias e processos estocásticos.  Rio de Janeiro, RJ: Interciência, 2008.  334p.  ISBN: 9788571931909. 
\\ 
 & \href{www.biblioteca.ifce.edu.br/index.asp?codigo_sophia=24604}{Meyer}, P. L.  Probabilidade: aplicações à estatística.  2º ed.  Rio de Janeiro, RJ : LTC, 1983.  426p.  ISBN: 8521602944.
\\ 
 &   \href{www.biblioteca.ifce.edu.br/index.asp?codigo_sophia=63227}{Navidi}, William.  Probabilidade e Estatística para Ciências Exatas.  1º ed.  Porto Alegre, RS : AMGH, 2012.  604p.  ISBN: 9788580550733.
\\
 &  \href{www.biblioteca.ifce.edu.br/index.asp?codigo_sophia=56742}{Bonafini}, Fernanda Cesar.  Estatística.  E-book.  Pearson.  186p.  ISBN: 9788564574403.
\\
 &  \href{www.biblioteca.ifce.edu.br/index.asp?codigo_sophia=55744}{Larson}, Ron; Farber, Elizabeth.  Estatística Aplicada.  E-book.  4º ed.  Pearson.  658p.  ISBN: 9788576053729. }

\def\PUDautor{Samuel Vieira Dias}

\def\PUDdata{2014-07-28}

\def\PUDrevisao{2}

\def\PUDrevisor{Samuel Vieira Dias}

\def\PUDdatarevisao{2017-05-09}

\PUDcorpo
